\chapter*{Introduction}
\addcontentsline{toc}{chapter}{Introduction}

Mon parcours à l'INSA Hauts-de-France en spécialité Mécatronique s'achève par le
stage de fin d'études, l'ultime étape avant le métier d'ingénieure. Après une
première expérience industrielle chez FORVIA en quatrième année, où j'avais
découvert la logistique, la conception mécanique et l'outillage, je souhaitais
pour ce PFE un projet qui rassemble tout ce qui m'a menée vers la mécatronique :
la robotique, l'embarqué, les réseaux et l'automobile.

C'est exactement ce que j'ai trouvé chez \textbf{ALTEN}, au sein des
\textbf{ALTEN Labs} de Sèvres, sur le projet \textbf{ProtoVA} : une plateforme
de véhicules autonomes à échelle réduite qui sert de démonstrateur aux
technologies du véhicule connecté. Mon stage, effectué du \fbox{JJ/MM/2026} au
\fbox{JJ/MM/2026}, avait pour ambition de transformer ces prototypes en
véhicules \emph{téléopérables}, \emph{autonomes}, \emph{communicants} et
\emph{supervisables}.

% TODO : ajuster les dates, puis relire — l'intro doit annoncer le plan.
Ce rapport s'organise comme suit : je présenterai d'abord ALTEN et
l'environnement dans lequel j'ai évolué, puis je dresserai un état de l'art des
domaines abordés. Après avoir posé la problématique et la démarche, je
détaillerai mes contributions en suivant la montée en puissance du véhicule :
fiabilisation de la plateforme mécanique (Protova1), téléopération WiFi et 5G,
perception et navigation autonome, communication véhicule-à-véhicule, et enfin
supervision temps réel. Un bilan global et les perspectives clôtureront ce
travail.
